\documentclass[UTF8,11pt,oneside,openany,fontset=none]{ctexbook}
\usepackage[a4paper,margin=23mm,top=25mm,bottom=25mm]{geometry}
\usepackage{fontspec}
\usepackage{amsmath,amssymb,enumitem,fancyhdr}
\usepackage[colorlinks=true,linkcolor=black,urlcolor=blue]{hyperref}
\setCJKmainfont{Noto Serif CJK SC}
\setCJKsansfont{Noto Sans CJK SC}
\setmainfont{DejaVu Sans}
\setlength{\parindent}{2em}
\setlength{\parskip}{0.4em}
\setlength{\headheight}{16pt}
\setlist[enumerate]{itemsep=0.3em,topsep=0.3em}
\pagestyle{fancy}
\fancyhf{}
\fancyhead[L]{化学教材来源独立数值质检}
\fancyhead[R]{研究隔离 · 未经授权}
\fancyfoot[C]{\thepage}
\sloppy
\title{大学化学数值复核质检册\\[0.4em]\large 42 道来源固定的科学核验记录}
\author{三千弱水 · 化学题库研究记录}
\date{2026 年 10 月}
\begin{document}
\frontmatter
\maketitle
\chapter*{版权及科学审核声明}
本册仅供内部学术核验，不复制原教材习题全文，不属于已授权中文翻译，也不提供可直接出版的题库。所有原文请打开每题的源文件链接。
其中 42 道有独立数值复算；34 道有三步中文计算评论；8 道为复核摘要。来源作品版本、单题内容权利、专家科学审核均未完成，禁止公开发布或用于正式组卷。
检测到 3 项参考答案标签或题号问题，以及 1 项化学式疑似笔误，仅作为待审事实记录。
\tableofcontents
\mainmatter
\chapter{无机化学}
\section{化学式中碳元素质量}
\noindent\textbf{来源章节：} contents/m51011.md；序号 1。\par
\noindent\textbf{原文链接：} \url{https://github.com/philschatz/chemistry-book/blob/a6ac3918e61cbedf0589c5b43e748dc8a7edf5e2/contents/m51011.md}\par
\noindent\textbf{固定 ID：} \texttt{OSCHEM-1957d6352e0f0f7c239a78cc}\par
\subsection*{数值审查摘要}
计算各化学式中碳原子的总质量；用 C=12.01 u。
\noindent 计算值：\texttt{12.01000}\par
\noindent 计算值：\texttt{12.01000}\par
\noindent 计算值：\texttt{144.1200}\par
\noindent 计算值：\texttt{60.05000}\par
\noindent\textit{待审核：科学专家审校、版本授权、单题版权、翻译及编排。不可用于正式出卷。}
\medskip\hrule\medskip
\section{五种化合物相对式量}
\noindent\textbf{来源章节：} contents/m51011.md；序号 3。\par
\noindent\textbf{原文链接：} \url{https://github.com/philschatz/chemistry-book/blob/a6ac3918e61cbedf0589c5b43e748dc8a7edf5e2/contents/m51011.md}\par
\noindent\textbf{固定 ID：} \texttt{OSCHEM-2ff4a6140f201ba1ac612626}\par
\subsection*{数值审查摘要}
按元素原子量分别计算 P4、H2O、Ca(NO3)2、乙酸、蔗糖的式量。
\noindent 计算值：\texttt{123.8960}\par
\noindent 计算值：\texttt{18.01500}\par
\noindent 计算值：\texttt{164.0860}\par
\noindent 计算值：\texttt{60.05200}\par
\noindent 计算值：\texttt{342.2970}\par
\noindent\textit{待审核：科学专家审校、版本授权、单题版权、翻译及编排。不可用于正式出卷。}
\medskip\hrule\medskip
\section{比较三种化合物中的氧元素质量}
\noindent\textbf{来源章节：} contents/m51011.md；序号 9。\par
\noindent\textbf{原文链接：} \url{https://github.com/philschatz/chemistry-book/blob/a6ac3918e61cbedf0589c5b43e748dc8a7edf5e2/contents/m51011.md}\par
\noindent\textbf{固定 ID：} \texttt{OSCHEM-5e552a98f4e66c44a51cebf7}\par
\subsection*{独立中文推导}
\begin{enumerate}
\item 乙醇与水每分子各含一个氧原子，甲酸含两个，因此不能只比较各化合物的摩尔数。
\item 对三种样品用 m(O)=n(化合物)×每粒子氧原子数×15.999 g/mol，分别取 0.75、0.60、1.00 mol。
\item 氧质量约为 12.00、19.20、16.00 g，故甲酸所含氧元素质量最多。
\end{enumerate}
\subsection*{结果与单位}
\noindent 乙醇中的氧：\texttt{11.99925} g\par
\noindent 甲酸中的氧：\texttt{19.19880} g\par
\noindent 水中的氧：\texttt{15.99900} g\par
\noindent\textit{待审核：科学专家审校、版本授权、单题版权、翻译及编排。不可用于正式出卷。}
\medskip\hrule\medskip
\section{五种物质的摩尔质量核算}
\noindent\textbf{来源章节：} contents/m51011.md；序号 13。\par
\noindent\textbf{原文链接：} \url{https://github.com/philschatz/chemistry-book/blob/a6ac3918e61cbedf0589c5b43e748dc8a7edf5e2/contents/m51011.md}\par
\noindent\textbf{固定 ID：} \texttt{OSCHEM-de6b22805e6871a63f7c759e}\par
\subsection*{独立中文推导}
\begin{enumerate}
\item 拆分分子式下标，尤其 Sc₂(SO₄)₃ 中 2 个 Sc、3 个 S、12 个 O；丙酮含三个 C。
\item 按 M=Σ n\_i A\_i 求和，例如 M(Sc₂(SO₄)₃)=2A(Sc)+3A(S)+12A(O)。
\item 所用标准原子量小数位可能与教材旧表不同，复核时先检查组成和量纲，再比较末位舍入。
\end{enumerate}
\subsection*{结果与单位}
\noindent S₈：\texttt{256.5280} g/mol\par
\noindent 戊烷：\texttt{72.15100} g/mol\par
\noindent 硫酸钪：\texttt{378.0800} g/mol\par
\noindent 丙酮：\texttt{58.08000} g/mol\par
\noindent 葡萄糖：\texttt{180.1560} g/mol\par
\noindent\textit{待审核：科学专家审校、版本授权、单题版权、翻译及编排。不可用于正式出卷。}
\medskip\hrule\medskip
\section{五种含杂原子有机物的式量}
\noindent\textbf{来源章节：} contents/m51011.md；序号 15。\par
\noindent\textbf{原文链接：} \url{https://github.com/philschatz/chemistry-book/blob/a6ac3918e61cbedf0589c5b43e748dc8a7edf5e2/contents/m51011.md}\par
\noindent\textbf{固定 ID：} \texttt{OSCHEM-7d8a3280e55269d3dd0acc3d}\par
\subsection*{独立中文推导}
\begin{enumerate}
\item 按已给物质的化学式计算而非只凭名称猜测盐型；含 Br、Cl、F 的组成尤其要逐原子计数。
\item 将 C₂HBrClF₃、C₁₂H₁₄N₂Cl₂、C₈H₁₀N₄O₂、CO(NH₂)₂、C₁₈H₃₅O₂Na 的各元素贡献求和。
\item 保留化学式和所用原子量表，源答案末位与现代原子量常有微差，须人工审定而非强行取齐。
\end{enumerate}
\subsection*{结果与单位}
\noindent 氟烷：\texttt{197.3780} g/mol\par
\noindent 百草枯盐式：\texttt{257.1580} g/mol\par
\noindent 咖啡因：\texttt{194.1940} g/mol\par
\noindent 尿素：\texttt{60.05600} g/mol\par
\noindent 钠皂：\texttt{306.4660} g/mol\par
\noindent\textit{待审核：科学专家审校、版本授权、单题版权、翻译及编排。不可用于正式出卷。}
\medskip\hrule\medskip
\section{等摩尔氧原子、氧气、臭氧的质量}
\noindent\textbf{来源章节：} contents/m51011.md；序号 21。\par
\noindent\textbf{原文链接：} \url{https://github.com/philschatz/chemistry-book/blob/a6ac3918e61cbedf0589c5b43e748dc8a7edf5e2/contents/m51011.md}\par
\noindent\textbf{固定 ID：} \texttt{OSCHEM-f765a89be00a4cd76cbb7263}\par
\subsection*{独立中文推导}
\begin{enumerate}
\item 三组都是 0.600 mol，但作为粒子的 O、O₂、O₃ 每粒子分别含 1、2、3 个氧原子。
\item 利用 m=nM，分别求 0.600×15.999、0.600×31.998、0.600×47.997 g。
\item 结果约 9.60、19.2、28.8 g；等摩尔粒子并不是等摩尔氧原子。
\end{enumerate}
\subsection*{结果与单位}
\noindent 氧原子：\texttt{9.599400} g\par
\noindent 氧气：\texttt{19.19880} g\par
\noindent 臭氧：\texttt{28.79820} g\par
\noindent\textit{待审核：科学专家审校、版本授权、单题版权、翻译及编排。不可用于正式出卷。}
\medskip\hrule\medskip
\section{锆石的元素质量与原子数}
\noindent\textbf{来源章节：} contents/m51011.md；序号 23。\par
\noindent\textbf{原文链接：} \url{https://github.com/philschatz/chemistry-book/blob/a6ac3918e61cbedf0589c5b43e748dc8a7edf5e2/contents/m51011.md}\par
\noindent\textbf{固定 ID：} \texttt{OSCHEM-dfc583d0b96b0b00c6c9d2ed}\par
\subsection*{独立中文推导}
\begin{enumerate}
\item 硅酸锆化学式 ZrSiO₄，0.3384 mol 中锆、硅各有 0.3384 mol 原子，而氧原子为四倍。
\item 原子数用 6.02214076×10²³ mol⁻¹ 求得；元素质量为各元素原子物质的量乘其摩尔质量。
\item 氧原子数必须是锆原子数四倍，结果分别标注“个”和“克”，避免混用元素数量与质量。
\end{enumerate}
\subsection*{结果与单位}
\noindent 锆原子数：\texttt{2.037892e+23} 个\par
\noindent 锆质量：\texttt{30.87020} g\par
\noindent 硅原子数：\texttt{2.037892e+23} 个\par
\noindent 硅质量：\texttt{9.503964} g\par
\noindent 氧原子数：\texttt{8.151570e+23} 个\par
\noindent 氧质量：\texttt{21.65625} g\par
\noindent\textit{待审核：科学专家审校、版本授权、单题版权、翻译及编排。不可用于正式出卷。}
\medskip\hrule\medskip
\section{三种铝盐中的铝质量比较}
\noindent\textbf{来源章节：} contents/m51011.md；序号 25。\par
\noindent\textbf{原文链接：} \url{https://github.com/philschatz/chemistry-book/blob/a6ac3918e61cbedf0589c5b43e748dc8a7edf5e2/contents/m51011.md}\par
\noindent\textbf{固定 ID：} \texttt{OSCHEM-c91d4c551c217cf7e5f5a7e2}\par
\textbf{化学式原文疑点：} printed\_Al2C1\_6\_is\_Al2Cl6\_not\_carbon\_digit\_one\par
\subsection*{独立中文推导}
\begin{enumerate}
\item 原题出现 C1 与 Cl 易混的化学式，当前只按原答案暗示的 Al₂Cl₆ 解释，保留原文疑点。
\item 对 122 g AlPO₄、266 g Al₂Cl₆、225 g Al₂S₃ 分别计算 m(Al)=m(样品)×n\_Al A(Al)/M(物质)。
\item 三者比较以 Al₂S₃ 含铝质量最大；若原化学式解释改变必须重算，不能自动抹除错字警告。
\end{enumerate}
\subsection*{结果与单位}
\noindent 磷酸铝中的 Al：\texttt{26.99223} g\par
\noindent 氯化铝二聚体中的 Al：\texttt{53.82883} g\par
\noindent 硫化铝中的 Al：\texttt{80.86741} g\par
\noindent\textit{待审核：科学专家审校、版本授权、单题版权、翻译及编排。不可用于正式出卷。}
\medskip\hrule\medskip
\section{钻石克拉质量换算至原子数}
\noindent\textbf{来源章节：} contents/m51011.md；序号 26。\par
\noindent\textbf{原文链接：} \url{https://github.com/philschatz/chemistry-book/blob/a6ac3918e61cbedf0589c5b43e748dc8a7edf5e2/contents/m51011.md}\par
\noindent\textbf{固定 ID：} \texttt{OSCHEM-fb0c34b160953464c4efad67}\par
\textbf{无教材参考答案：} 数值已独立求取，尚待人工复核。\par
\subsection*{独立中文推导}
\begin{enumerate}
\item 按题设 1 克拉等于 0.200 g，1.25 克拉钻石质量就是 0.250 g。
\item 假设金刚石纯碳，先计算 n=0.250/12.011 mol，再乘阿伏伽德罗常数 N\_A。
\item 得到约 1.253×10²² 个原子；源教材无附答案，钻石纯度是应单独审阅的假设。
\end{enumerate}
\subsection*{结果与单位}
\noindent 碳原子数：\texttt{1.253464e+22} 个\par
\noindent\textit{待审核：科学专家审校、版本授权、单题版权、翻译及编排。不可用于正式出卷。}
\medskip\hrule\medskip
\chapter{分析化学}
\section{溶液摩尔浓度核验}
\noindent\textbf{来源章节：} contents/m51014.md；序号 4。\par
\noindent\textbf{原文链接：} \url{https://github.com/philschatz/chemistry-book/blob/a6ac3918e61cbedf0589c5b43e748dc8a7edf5e2/contents/m51014.md}\par
\noindent\textbf{固定 ID：} \texttt{OSCHEM-6353a589542e63d752c60053}\par
\subsection*{数值审查摘要}
质量先除摩尔质量，毫升换算为升，水合物按完整式量；原题含 MathML，正式渲染仍待审核。
\noindent 计算值：\texttt{0.6788991}\par
\noindent 计算值：\texttt{1.000061}\par
\noindent 计算值：\texttt{0.06998056}\par
\noindent 计算值：\texttt{1.752144}\par
\noindent 计算值：\texttt{0.07000000}\par
\noindent 计算值：\texttt{6.582917}\par
\noindent\textit{待审核：科学专家审校、版本授权、单题版权、翻译及编排。不可用于正式出卷。}
\medskip\hrule\medskip
\section{葡萄糖定量配液}
\noindent\textbf{来源章节：} contents/m51014.md；序号 6。\par
\noindent\textbf{原文链接：} \url{https://github.com/philschatz/chemistry-book/blob/a6ac3918e61cbedf0589c5b43e748dc8a7edf5e2/contents/m51014.md}\par
\noindent\textbf{固定 ID：} \texttt{OSCHEM-28cc53aec9bf5a34cad7ccdc}\par
\subsection*{数值审查摘要}
糖的物质的量为 0.500×0.30 mol，质量为 0.15×180.156 g。
\noindent 计算值：\texttt{27.02340}\par
\noindent\textit{待审核：科学专家审校、版本授权、单题版权、翻译及编排。不可用于正式出卷。}
\medskip\hrule\medskip
\section{高锰酸钾质量配液浓度}
\noindent\textbf{来源章节：} contents/m51014.md；序号 10。\par
\noindent\textbf{原文链接：} \url{https://github.com/philschatz/chemistry-book/blob/a6ac3918e61cbedf0589c5b43e748dc8a7edf5e2/contents/m51014.md}\par
\noindent\textbf{固定 ID：} \texttt{OSCHEM-94a2f4f4ce444c31880fc9ac}\par
\subsection*{独立中文推导}
\begin{enumerate}
\item KMnO₄ 摩尔质量等于 A(K)+A(Mn)+4A(O)，在计算前逐一清点下标。
\item 先用 0.0908 g 除以式量得到摩尔数，再除以 0.500 L 得浓度 c。
\item 约 1.149×10⁻³ mol/L，有效数字约 1.15×10⁻³ mol/L；逆算 cVM 应回到称样质量。
\end{enumerate}
\subsection*{结果与单位}
\noindent KMnO₄ 浓度：\texttt{0.001149132} mol/L\par
\noindent\textit{待审核：科学专家审校、版本授权、单题版权、翻译及编排。不可用于正式出卷。}
\medskip\hrule\medskip
\section{胆固醇、氨、异丙醇和碘的浓度}
\noindent\textbf{来源章节：} contents/m51014.md；序号 12。\par
\noindent\textbf{原文链接：} \url{https://github.com/philschatz/chemistry-book/blob/a6ac3918e61cbedf0589c5b43e748dc8a7edf5e2/contents/m51014.md}\par
\noindent\textbf{固定 ID：} \texttt{OSCHEM-b692102969d51c51ea7c870f}\par
\subsection*{独立中文推导}
\begin{enumerate}
\item 四组都先检查溶质化学式和质量单位，异丙醇质量 1.49 kg 须转换为 1490 g。
\item 分别计算 c=m/(MV)，体积依次 0.100、0.500、2.50、0.100 L；摩尔质量分别按 C₂₇H₄₆O、NH₃、C₃H₈O、I₂。
\item 碘为双原子分子，误用单个碘原子量会造成两倍偏差；其他数据也要按各自精度舍入。
\end{enumerate}
\subsection*{结果与单位}
\noindent 胆固醇：\texttt{0.005043138} mol/L\par
\noindent 氨：\texttt{0.4990899} mol/L\par
\noindent 异丙醇：\texttt{9.917465} mol/L\par
\noindent 碘：\texttt{0.001142592} mol/L\par
\noindent\textit{待审核：科学专家审校、版本授权、单题版权、翻译及编排。不可用于正式出卷。}
\medskip\hrule\medskip
\section{牛奶中钙离子的浓度}
\noindent\textbf{来源章节：} contents/m51014.md；序号 14。\par
\noindent\textbf{原文链接：} \url{https://github.com/philschatz/chemistry-book/blob/a6ac3918e61cbedf0589c5b43e748dc8a7edf5e2/contents/m51014.md}\par
\noindent\textbf{固定 ID：} \texttt{OSCHEM-d7533f0e2adbcdaf08005e6d}\par
\subsection*{独立中文推导}
\begin{enumerate}
\item 题设含 1.0 g 钙元素（作为 Ca²⁺），不是 1.0 g 的其他钙盐。
\item n(Ca²⁺)=1.0 g/(40.078 g/mol)，再除以 1.0 L 溶液体积。
\item 约 0.02495 mol/L，按输入有效数字取 0.025 mol/L，需区分元素质量与化合物质量。
\end{enumerate}
\subsection*{结果与单位}
\noindent 钙离子浓度：\texttt{0.02495134} mol/L\par
\noindent\textit{待审核：科学专家审校、版本授权、单题版权、翻译及编排。不可用于正式出卷。}
\medskip\hrule\medskip
\section{硝酸铁的原液取样体积}
\noindent\textbf{来源章节：} contents/m51014.md；序号 15。\par
\noindent\textbf{原文链接：} \url{https://github.com/philschatz/chemistry-book/blob/a6ac3918e61cbedf0589c5b43e748dc8a7edf5e2/contents/m51014.md}\par
\noindent\textbf{固定 ID：} \texttt{OSCHEM-56ce2afd36d22a25143d372a}\par
\textbf{无教材参考答案：} 数值已独立求取，尚待人工复核。\par
\subsection*{独立中文推导}
\begin{enumerate}
\item 原液 1.00 mol/L，目标为 1.00 L、0.250 mol/L；假设稀释前后溶质的量守恒。
\item 依 c₁V₁=c₂V₂，V₁=(0.250×1.00)/1.00 L。
\item 计算 0.250 L 即 250 mL；取样体积小于最终溶液体积，是合理性检查。
\end{enumerate}
\subsection*{结果与单位}
\noindent 原液体积：\texttt{0.2500000} L\par
\noindent\textit{待审核：科学专家审校、版本授权、单题版权、翻译及编排。不可用于正式出卷。}
\medskip\hrule\medskip
\section{异丙醇稀释后的总体积}
\noindent\textbf{来源章节：} contents/m51014.md；序号 16。\par
\noindent\textbf{原文链接：} \url{https://github.com/philschatz/chemistry-book/blob/a6ac3918e61cbedf0589c5b43e748dc8a7edf5e2/contents/m51014.md}\par
\noindent\textbf{固定 ID：} \texttt{OSCHEM-150c01a8a661e6084d0e5d12}\par
\subsection*{独立中文推导}
\begin{enumerate}
\item 已知稀释只加溶剂，溶质的物质的量在操作前后不变；体积、浓度必须使用配套单位。
\item 写出 c₁V₁=c₂V₂，并将 0.3556 mol/L、0.1718 L 与 0.1222 mol/L 代入反求 V₂。
\item V₂≈0.499935 L，有效数字约 0.5000 L；用 V₂c₂=V₁c₁ 做反向一致性检查。
\end{enumerate}
\subsection*{结果与单位}
\noindent 最终体积：\texttt{0.4999352} L\par
\noindent\textit{待审核：科学专家审校、版本授权、单题版权、翻译及编排。不可用于正式出卷。}
\medskip\hrule\medskip
\section{蔗糖储备液取用体积}
\noindent\textbf{来源章节：} contents/m51014.md；序号 18。\par
\noindent\textbf{原文链接：} \url{https://github.com/philschatz/chemistry-book/blob/a6ac3918e61cbedf0589c5b43e748dc8a7edf5e2/contents/m51014.md}\par
\noindent\textbf{固定 ID：} \texttt{OSCHEM-05dc2a3a7b98e9cab288a160}\par
\subsection*{独立中文推导}
\begin{enumerate}
\item 储备液浓度 0.33 mol/L 大于最终 0.025 mol/L，原液取样体积应小于总体积 25 mL。
\item 稀释不改变溶质物质的量，V₁=c₂V₂/c₁=(0.025×25)/0.33 mL。
\item 结果约 1.89394 mL，按两位有效数字约 1.9 mL；以原液溶质的量反查。
\end{enumerate}
\subsection*{结果与单位}
\noindent 原液体积：\texttt{1.893939} mL\par
\noindent\textit{待审核：科学专家审校、版本授权、单题版权、翻译及编排。不可用于正式出卷。}
\medskip\hrule\medskip
\section{四组溶液稀释及原答案题号勘误}
\noindent\textbf{来源章节：} contents/m51014.md；序号 20。\par
\noindent\textbf{原文链接：} \url{https://github.com/philschatz/chemistry-book/blob/a6ac3918e61cbedf0589c5b43e748dc8a7edf5e2/contents/m51014.md}\par
\noindent\textbf{固定 ID：} \texttt{OSCHEM-bd8233f4402178ae790b95d8}\par
\textbf{参考答案差异：} answer\_labels\_part\_d\_as\_part\_e\par
\subsection*{独立中文推导}
\begin{enumerate}
\item 所有小问用 c\_f=c\_iV\_i/V\_f，确保体积单位在每个小问中前后一致。
\item 四个浓度计算分别为 (1.00×0.250)/2.00、(0.5000×0.1222)/1.250、(2.35×0.350)/4.00、(22.50×0.025)/100.0。
\item 第四个小问教材参考答案误写成 (e) 而非 (d)，这属于编号错位，不影响溶质守恒的公式。
\end{enumerate}
\subsection*{结果与单位}
\noindent Fe(NO₃)₃：\texttt{0.1250000} mol/L\par
\noindent 丙醇：\texttt{0.04888000} mol/L\par
\noindent H₃PO₄：\texttt{0.2056250} mol/L\par
\noindent 蔗糖：\texttt{0.005625000} mol/L\par
\noindent\textit{待审核：科学专家审校、版本授权、单题版权、翻译及编排。不可用于正式出卷。}
\medskip\hrule\medskip
\section{HCl 储备液的摩尔浓度}
\noindent\textbf{来源章节：} contents/m51014.md；序号 22。\par
\noindent\textbf{原文链接：} \url{https://github.com/philschatz/chemistry-book/blob/a6ac3918e61cbedf0589c5b43e748dc8a7edf5e2/contents/m51014.md}\par
\noindent\textbf{固定 ID：} \texttt{OSCHEM-21ce27f2ffa9de7f66f4308f}\par
\subsection*{独立中文推导}
\begin{enumerate}
\item 明确 868.8 g 是 HCl 溶质质量，2.00 L 是溶液定容体积，不可直接把 g/L 视为 mol/L。
\item 计算 M(HCl)=1.008+35.45=36.458 g/mol，由 c=m/(MV) 得 868.8/(36.458×2.00)。
\item c≈11.9144 mol/L，按原始量精度约 11.9 mol/L；复算 cVM 应为 868.8 g。
\end{enumerate}
\subsection*{结果与单位}
\noindent HCl 浓度：\texttt{11.91508} mol/L\par
\noindent\textit{待审核：科学专家审校、版本授权、单题版权、翻译及编排。不可用于正式出卷。}
\medskip\hrule\medskip
\section{硫酸钾溶液的总体积}
\noindent\textbf{来源章节：} contents/m51014.md；序号 24。\par
\noindent\textbf{原文链接：} \url{https://github.com/philschatz/chemistry-book/blob/a6ac3918e61cbedf0589c5b43e748dc8a7edf5e2/contents/m51014.md}\par
\noindent\textbf{固定 ID：} \texttt{OSCHEM-9fd08976c31cc12804d2a9b4}\par
\subsection*{独立中文推导}
\begin{enumerate}
\item 化学式 K₂SO₄ 含两 K、一 S、四 O，先求摩尔质量，防止漏乘下标。
\item M=2×39.0983+32.06+4×15.999 g/mol；由 n=m/M、V=n/c 得 V=(57/M)/0.20。
\item 求得约 1.63556 L，按给定精度约 1.6 L；由 cVM 验算溶质质量。
\end{enumerate}
\subsection*{结果与单位}
\noindent 体积：\texttt{1.635557} L\par
\noindent\textit{待审核：科学专家审校、版本授权、单题版权、翻译及编排。不可用于正式出卷。}
\medskip\hrule\medskip
\section{40 ℃ 下纯水的浓度与 pH}
\noindent\textbf{来源章节：} contents/m51117.md；序号 2。\par
\noindent\textbf{原文链接：} \url{https://github.com/philschatz/chemistry-book/blob/a6ac3918e61cbedf0589c5b43e748dc8a7edf5e2/contents/m51117.md}\par
\noindent\textbf{固定 ID：} \texttt{OSCHEM-39061d8176054530dd51bfdd}\par
\textbf{无教材参考答案：} 数值已独立求取，尚待人工复核。\par
\subsection*{独立中文推导}
\begin{enumerate}
\item 纯水中性条件是 [H₃O⁺]=[OH⁻]，而不是任何温度下 pH 均为 7。
\item 40 ℃ 时 K\_w=2.9×10⁻¹⁴，因此 [H₃O⁺]=√K\_w≈1.70294×10⁻⁷ mol/L。
\item pH=−log₁₀([H₃O⁺]/1 mol/L)≈6.7688；教材无参考答案，活度近似应另行复核。
\end{enumerate}
\subsection*{结果与单位}
\noindent 水合氢离子浓度：\texttt{1.702939e-7} mol/L\par
\noindent pH：\texttt{6.768801} \par
\noindent\textit{待审核：科学专家审校、版本授权、单题版权、翻译及编排。不可用于正式出卷。}
\medskip\hrule\medskip
\section{纯水非室温酸碱平衡}
\noindent\textbf{来源章节：} contents/m51117.md；序号 3。\par
\noindent\textbf{原文链接：} \url{https://github.com/philschatz/chemistry-book/blob/a6ac3918e61cbedf0589c5b43e748dc8a7edf5e2/contents/m51117.md}\par
\noindent\textbf{固定 ID：} \texttt{OSCHEM-8b257cdd67b8a3aeb88bb6e3}\par
\subsection*{数值审查摘要}
纯水中 [H3O+]=[OH−]=sqrt(Kw)，pH=pOH=−log10(sqrt(Kw))。
\noindent 计算值：\texttt{3.051393e-7}\par
\noindent 计算值：\texttt{6.515502}\par
\noindent 计算值：\texttt{6.515502}\par
\noindent\textit{待审核：科学专家审校、版本授权、单题版权、翻译及编排。不可用于正式出卷。}
\medskip\hrule\medskip
\section{强酸强碱 pH 与 pOH 勘误}
\noindent\textbf{来源章节：} contents/m51117.md；序号 5。\par
\noindent\textbf{原文链接：} \url{https://github.com/philschatz/chemistry-book/blob/a6ac3918e61cbedf0589c5b43e748dc8a7edf5e2/contents/m51117.md}\par
\noindent\textbf{固定 ID：} \texttt{OSCHEM-5d36939cdd0bbd60e8e95441}\par
\textbf{参考答案差异：} textbook\_swaps\_pH\_and\_pOH\_for\_b\_and\_d\par
\subsection*{数值审查摘要}
按题目完全电离及理想浓度近似，pH+pOH=14。教材答案将 (b) NaOH 与 (d) KOH 的 pH、pOH 名称互换，必须指出，浓溶液还需考虑活度。
\noindent 计算值：\texttt{3.586700}\par
\noindent 计算值：\texttt{10.41330}\par
\noindent 计算值：\texttt{13.32222}\par
\noindent 计算值：\texttt{0.6777807}\par
\noindent 计算值：\texttt{10.15229}\par
\noindent 计算值：\texttt{3.847712}\par
\noindent 计算值：\texttt{14.39794}\par
\noindent 计算值：\texttt{-0.3979400}\par
\noindent\textit{待审核：科学专家审校、版本授权、单题版权、翻译及编排。不可用于正式出卷。}
\medskip\hrule\medskip
\section{从 pH 独立求两种离子浓度}
\noindent\textbf{来源章节：} contents/m51117.md；序号 7。\par
\noindent\textbf{原文链接：} \url{https://github.com/philschatz/chemistry-book/blob/a6ac3918e61cbedf0589c5b43e748dc8a7edf5e2/contents/m51117.md}\par
\noindent\textbf{固定 ID：} \texttt{OSCHEM-288afabff92ecbad81ec8e97}\par
\subsection*{独立中文推导}
\begin{enumerate}
\item 稀溶液下近似以浓度代替活度，25 ℃ 时采用 K\_w≈10⁻¹⁴。
\item 由 pH=6.52 求 [H₃O⁺]≈10⁻⁶·⁵²，再以 [OH⁻]≈K\_w/[H₃O⁺] 计算。
\item 两者分别约 3.020×10⁻⁷ 和 3.311×10⁻⁸ mol/L，乘积反算约 10⁻¹⁴。
\end{enumerate}
\subsection*{结果与单位}
\noindent 水合氢离子浓度：\texttt{3.019952e-7} mol/L\par
\noindent 氢氧根离子浓度：\texttt{3.311311e-8} mol/L\par
\noindent\textit{待审核：科学专家审校、版本授权、单题版权、翻译及编排。不可用于正式出卷。}
\medskip\hrule\medskip
\section{雨水中氢氧根离子浓度}
\noindent\textbf{来源章节：} contents/m51117.md；序号 10。\par
\noindent\textbf{原文链接：} \url{https://github.com/philschatz/chemistry-book/blob/a6ac3918e61cbedf0589c5b43e748dc8a7edf5e2/contents/m51117.md}\par
\noindent\textbf{固定 ID：} \texttt{OSCHEM-7538ab68d00552e9a91ef013}\par
\textbf{无教材参考答案：} 数值已独立求取，尚待人工复核。\par
\subsection*{独立中文推导}
\begin{enumerate}
\item 25 ℃ 时按稀溶液近似，取 K\_w≈1.0×10⁻¹⁴。
\item 雨水 [H₃O⁺]=1.7×10⁻⁶ mol/L，代入 [OH⁻]=K\_w/[H₃O⁺]。
\item 得到 [OH⁻]≈5.882×10⁻⁹ mol/L；两浓度相乘反验 K\_w，教材原题未附答案。
\end{enumerate}
\subsection*{结果与单位}
\noindent 氢氧根浓度：\texttt{5.882353e-9} mol/L\par
\noindent\textit{待审核：科学专家审校、版本授权、单题版权、翻译及编排。不可用于正式出卷。}
\medskip\hrule\medskip
\section{氨水离子积及原参考答案标签错误}
\noindent\textbf{来源章节：} contents/m51117.md；序号 11。\par
\noindent\textbf{原文链接：} \url{https://github.com/philschatz/chemistry-book/blob/a6ac3918e61cbedf0589c5b43e748dc8a7edf5e2/contents/m51117.md}\par
\noindent\textbf{固定 ID：} \texttt{OSCHEM-aa57c9f57f03106d8295119b}\par
\textbf{参考答案差异：} textbook\_labels\_hydronium\_numeric\_result\_as\_hydroxide\par
\subsection*{独立中文推导}
\begin{enumerate}
\item 原数据给定的是 [OH⁻]=3.2×10⁻³ mol/L，但要求计算的是 [H₃O⁺]，必须辨别符号。
\item 25 ℃ 下按 [H₃O⁺]=K\_w/[OH⁻]，代入 10⁻¹⁴/(3.2×10⁻³)。
\item 结果约 3.125×10⁻¹² mol/L；原答案将该数值误标为 OH⁻，必须作为来源疑点保留。
\end{enumerate}
\subsection*{结果与单位}
\noindent 水合氢离子浓度：\texttt{3.125000e-12} mol/L\par
\noindent\textit{待审核：科学专家审校、版本授权、单题版权、翻译及编排。不可用于正式出卷。}
\medskip\hrule\medskip
\chapter{物理化学}
\section{冷热水混合量热}
\noindent\textbf{来源章节：} contents/m51034.md；序号 5。\par
\noindent\textbf{原文链接：} \url{https://github.com/philschatz/chemistry-book/blob/a6ac3918e61cbedf0589c5b43e748dc8a7edf5e2/contents/m51034.md}\par
\noindent\textbf{固定 ID：} \texttt{OSCHEM-189c66c8ab6a1ee34fa1c983}\par
\textbf{无教材参考答案：} 数值已独立求取，尚待人工复核。\par
\subsection*{数值审查摘要}
在同密度和同热容、无热损假设下，以 m\_hot(95−60)=m\_cold(60−23) 求冷水体积。
\noindent 计算值：\texttt{170.2703}\par
\noindent\textit{待审核：科学专家审校、版本授权、单题版权、翻译及编排。不可用于正式出卷。}
\medskip\hrule\medskip
\section{热咖啡与银勺的能量守恒}
\noindent\textbf{来源章节：} contents/m51034.md；序号 6。\par
\noindent\textbf{原文链接：} \url{https://github.com/philschatz/chemistry-book/blob/a6ac3918e61cbedf0589c5b43e748dc8a7edf5e2/contents/m51034.md}\par
\noindent\textbf{固定 ID：} \texttt{OSCHEM-0e809cbd3169cbd6ba3864a6}\par
\subsection*{数值审查摘要}
由咖啡放热和银勺吸热守恒计算咖啡温降，使用 c\_water=4.184 J/(g·K)。
\noindent 计算值：\texttt{0.9896324}\par
\noindent\textit{待审核：科学专家审校、版本授权、单题版权、翻译及编排。不可用于正式出卷。}
\medskip\hrule\medskip
\section{汽车冷却水降温传热}
\noindent\textbf{来源章节：} contents/m51034.md；序号 8。\par
\noindent\textbf{原文链接：} \url{https://github.com/philschatz/chemistry-book/blob/a6ac3918e61cbedf0589c5b43e748dc8a7edf5e2/contents/m51034.md}\par
\noindent\textbf{固定 ID：} \texttt{OSCHEM-45bcc141496551da3493422d}\par
\subsection*{独立中文推导}
\begin{enumerate}
\item 按美制 1 gallon≈3.78541 L、密度约 1 g/mL 得水质量约 3785 g。
\item 华氏温差 240−175=65 °F，相当于 65×5/9 K；因此 Q=mcΔT。
\item 计算热量绝对值约 572 kJ，按题设精度为 5.7×10² kJ；系统取水时放热符号为负。
\end{enumerate}
\subsection*{结果与单位}
\noindent 传热量绝对值：\texttt{571.9337} kJ\par
\noindent\textit{待审核：科学专家审校、版本授权、单题版权、翻译及编排。不可用于正式出卷。}
\medskip\hrule\medskip
\section{吸热反应的近似量热}
\noindent\textbf{来源章节：} contents/m51034.md；序号 14。\par
\noindent\textbf{原文链接：} \url{https://github.com/philschatz/chemistry-book/blob/a6ac3918e61cbedf0589c5b43e748dc8a7edf5e2/contents/m51034.md}\par
\noindent\textbf{固定 ID：} \texttt{OSCHEM-4d1c0cfafbad4476744d77b0}\par
\subsection*{独立中文推导}
\begin{enumerate}
\item 温度下降 3.1 ℃，说明反应体系吸热；近似按溶液质量与水相近的比热处理。
\item m≈100+3.15+1.52 g，Q≈m×4.20×3.1 J，再除以 1000。
\item 约 1.363 kJ（教材约 1.4 kJ）；真实值须考虑量热器热容和溶液比热偏差。
\end{enumerate}
\subsection*{结果与单位}
\noindent 吸热量：\texttt{1.362803} kJ\par
\noindent\textit{待审核：科学专家审校、版本授权、单题版权、翻译及编排。不可用于正式出卷。}
\medskip\hrule\medskip
\section{弹式量热热容与放热量}
\noindent\textbf{来源章节：} contents/m51034.md；序号 18。\par
\noindent\textbf{原文链接：} \url{https://github.com/philschatz/chemistry-book/blob/a6ac3918e61cbedf0589c5b43e748dc8a7edf5e2/contents/m51034.md}\par
\noindent\textbf{固定 ID：} \texttt{OSCHEM-6e4024a05f0ca869a990e0ba}\par
\subsection*{独立中文推导}
\begin{enumerate}
\item 量热器热容 534 J/K，675 mL 水在近似密度 1 g/mL 时热容为 675×4.184 J/K。
\item 总热容为两项相加，温差为 26.9−23.4=3.5 K；Q=C\_totalΔT。
\item 绝对值为 11.7537 kJ，教材舍入约 11.7 kJ；燃烧反应作为系统的热量符号为负。
\end{enumerate}
\subsection*{结果与单位}
\noindent 热量绝对值：\texttt{11.75370} kJ\par
\noindent\textit{待审核：科学专家审校、版本授权、单题版权、翻译及编排。不可用于正式出卷。}
\medskip\hrule\medskip
\section{脂肪在膳食中的能量占比}
\noindent\textbf{来源章节：} contents/m51034.md；序号 20。\par
\noindent\textbf{原文链接：} \url{https://github.com/philschatz/chemistry-book/blob/a6ac3918e61cbedf0589c5b43e748dc8a7edf5e2/contents/m51034.md}\par
\noindent\textbf{固定 ID：} \texttt{OSCHEM-02ce9d5faf259c98e2ca6374}\par
\subsection*{独立中文推导}
\begin{enumerate}
\item 食品 Cal 为千卡；脂肪 65 g、能量密度 9.1 Cal/g，已知每日总能量 2000 Cal。
\item 脂肪提供 65×9.1=591.5 Cal，故占比=591.5/2000×100\%。
\item 计算为 29.575\%，按题设精度取约 30\%；分子分母必须使用同一能量单位。
\end{enumerate}
\subsection*{结果与单位}
\noindent 能量百分比：\texttt{29.57500} \%\par
\noindent\textit{待审核：科学专家审校、版本授权、单题版权、翻译及编排。不可用于正式出卷。}
\medskip\hrule\medskip
\section{低热量饮料中的碳水质量严格上限}
\noindent\textbf{来源章节：} contents/m51034.md；序号 22。\par
\noindent\textbf{原文链接：} \url{https://github.com/philschatz/chemistry-book/blob/a6ac3918e61cbedf0589c5b43e748dc8a7edf5e2/contents/m51034.md}\par
\noindent\textbf{固定 ID：} \texttt{OSCHEM-5a5763b05e3aa1de5cac50cb}\par
\subsection*{独立中文推导}
\begin{enumerate}
\item 已给碳水化合物能量密度 4.1 Cal/g，若饮料热量严格小于 1 Cal，应该保留不等号。
\item 令质量为 m，列 4.1m<1，得 m<1/4.1 g。
\item 边界值约 0.243902 g，教材的 0.24 g 是近似说明；边界等号自身不符合严格小于条件。
\end{enumerate}
\subsection*{结果与单位}
\noindent 质量上限：\texttt{0.2439024} g\par
\noindent\textit{待审核：科学专家审校、版本授权、单题版权、翻译及编排。不可用于正式出卷。}
\medskip\hrule\medskip
\section{谷物食品的营养热量}
\noindent\textbf{来源章节：} contents/m51034.md；序号 24。\par
\noindent\textbf{原文链接：} \url{https://github.com/philschatz/chemistry-book/blob/a6ac3918e61cbedf0589c5b43e748dc8a7edf5e2/contents/m51034.md}\par
\noindent\textbf{固定 ID：} \texttt{OSCHEM-d3a5547d367f2eb13d3e0e96}\par
\subsection*{独立中文推导}
\begin{enumerate}
\item 食品单位 Cal 是千卡；蛋白质 3 g、碳水 18 g、脂肪 6 g。
\item 能量=3×4.1+18×4.1+6×9.1 Cal，以各营养素的单位质量能量分别求和。
\item 计算为 140.7 Cal，按给定精度约 1.4×10² Cal。
\end{enumerate}
\subsection*{结果与单位}
\noindent 每份热量：\texttt{140.7000} Cal\par
\noindent\textit{待审核：科学专家审校、版本授权、单题版权、翻译及编排。不可用于正式出卷。}
\medskip\hrule\medskip
\section{定容气体受热的理论压强}
\noindent\textbf{来源章节：} contents/m51065.md；序号 10。\par
\noindent\textbf{原文链接：} \url{https://github.com/philschatz/chemistry-book/blob/a6ac3918e61cbedf0589c5b43e748dc8a7edf5e2/contents/m51065.md}\par
\noindent\textbf{固定 ID：} \texttt{OSCHEM-ec63c02e19921375f30b30be}\par
\subsection*{独立中文推导}
\begin{enumerate}
\item 只在理想气体、定容和摩尔数不变的数学模型中使用 P/T=常数；摄氏温度不能直接相除。
\item 将 23 ℃ 与 475 ℃ 转为 296.15 K 与 748.15 K，代入 P₂=1344×748.15/296.15 torr。
\item 约 3395.28 torr，即 3.40×10³ torr；这是理论计算，不代表实际加热压力容器安全。
\end{enumerate}
\subsection*{结果与单位}
\noindent 理论末态压强：\texttt{3395.285} torr\par
\noindent\textit{待审核：科学专家审校、版本授权、单题版权、翻译及编排。不可用于正式出卷。}
\medskip\hrule\medskip
\section{等压氢气升温时的体积}
\noindent\textbf{来源章节：} contents/m51065.md；序号 12。\par
\noindent\textbf{原文链接：} \url{https://github.com/philschatz/chemistry-book/blob/a6ac3918e61cbedf0589c5b43e748dc8a7edf5e2/contents/m51065.md}\par
\noindent\textbf{固定 ID：} \texttt{OSCHEM-84bc4a73b1cbb198bae95007}\par
\subsection*{独立中文推导}
\begin{enumerate}
\item 假设气体的量和压强不变，运用 V/T=常数，先将 −196 ℃ 和 100 ℃ 转为绝对温度。
\item 两端温度分别是 77.15 K、373.15 K，因此 V₂=2.50×373.15/77.15 L。
\item 结果约 12.0917 L，按三位有效数字 12.1 L；以 V₂/T₂ 与 V₁/T₁ 的相等性反查。
\end{enumerate}
\subsection*{结果与单位}
\noindent 末态体积：\texttt{12.09170} L\par
\noindent\textit{待审核：科学专家审校、版本授权、单题版权、翻译及编排。不可用于正式出卷。}
\medskip\hrule\medskip
\section{三氟化硼的气体摩尔数与质量}
\noindent\textbf{来源章节：} contents/m51065.md；序号 16。\par
\noindent\textbf{原文链接：} \url{https://github.com/philschatz/chemistry-book/blob/a6ac3918e61cbedf0589c5b43e748dc8a7edf5e2/contents/m51065.md}\par
\noindent\textbf{固定 ID：} \texttt{OSCHEM-ea9385920480250198ff73fd}\par
\subsection*{独立中文推导}
\begin{enumerate}
\item 采用 P=1.220 atm、V=4.3410 L、T=788.0 K 和 R=0.082057338 L·atm/(mol·K)。
\item 先求 n=PV/(RT)，再以 M(BF₃)=10.81+3×18.998 g/mol 计算 m=nM。
\item 得到 n≈0.0819042 mol 与 m≈5.55343 g；m/M 反推 n 应与理想气体计算一致。
\end{enumerate}
\subsection*{结果与单位}
\noindent 物质的量：\texttt{0.08190416} mol\par
\noindent 质量：\texttt{5.553430} g\par
\noindent\textit{待审核：科学专家审校、版本授权、单题版权、翻译及编排。不可用于正式出卷。}
\medskip\hrule\medskip
\section{碘蒸气的理想气体温度}
\noindent\textbf{来源章节：} contents/m51065.md；序号 17。\par
\noindent\textbf{原文链接：} \url{https://github.com/philschatz/chemistry-book/blob/a6ac3918e61cbedf0589c5b43e748dc8a7edf5e2/contents/m51065.md}\par
\noindent\textbf{固定 ID：} \texttt{OSCHEM-abc4eca275a5634e93c948b3}\par
\textbf{无教材参考答案：} 数值已独立求取，尚待人工复核。\par
\subsection*{独立中文推导}
\begin{enumerate}
\item 碘蒸气分子 I₂ 的摩尔质量为 2×126.90447 g/mol，不能按单个 I 原子求摩尔数。
\item 气室 73.3 mL 转成 0.0733 L；n=0.292/M(I₂)，T=0.462×0.0733/(nR)。
\item 约为 358.72 K（约 85.6 ℃）。原题不附答案，气体模型还需人工审查。
\end{enumerate}
\subsection*{结果与单位}
\noindent 绝对温度：\texttt{358.7175} K\par
\noindent\textit{待审核：科学专家审校、版本授权、单题版权、翻译及编排。不可用于正式出卷。}
\medskip\hrule\medskip
\section{三种气体的压力体积单位换算}
\noindent\textbf{来源章节：} contents/m51065.md；序号 18。\par
\noindent\textbf{原文链接：} \url{https://github.com/philschatz/chemistry-book/blob/a6ac3918e61cbedf0589c5b43e748dc8a7edf5e2/contents/m51065.md}\par
\noindent\textbf{固定 ID：} \texttt{OSCHEM-56f52c350b8d228a760140c2}\par
\subsection*{独立中文推导}
\begin{enumerate}
\item 三份数据的压力分别使用 torr、kPa、torr；计算前必须统一为 atm，221 mL 改为 0.221 L。
\item 三组分别用 n=PV/(RT)，温度为 299.15、483、219.15 K，再乘各自气体摩尔质量。
\item 三份质量约 7.24×10⁻²、23.12、1.49×10⁻⁴ g；注意第一项是 0.100 L，不能错用 100 L。
\end{enumerate}
\subsection*{结果与单位}
\noindent CO₂：\texttt{0.07242022} g\par
\noindent C₂H₄：\texttt{23.12374} g\par
\noindent Ar：\texttt{0.0001485739} g\par
\noindent\textit{待审核：科学专家审校、版本授权、单题版权、翻译及编排。不可用于正式出卷。}
\medskip\hrule\medskip
\section{高压氧气折算到人体温度与常压}
\noindent\textbf{来源章节：} contents/m51065.md；序号 20。\par
\noindent\textbf{原文链接：} \url{https://github.com/philschatz/chemistry-book/blob/a6ac3918e61cbedf0589c5b43e748dc8a7edf5e2/contents/m51065.md}\par
\noindent\textbf{固定 ID：} \texttt{OSCHEM-08497f0802b2db9edb7127d5}\par
\subsection*{独立中文推导}
\begin{enumerate}
\item 初态为 151 atm、35.4 L、298.15 K，目标为 1 atm、310.15 K，假设理想气体和物质的量不变。
\item 使用 P₁V₁/T₁=P₂V₂/T₂，得到 V₂=151×35.4×310.15/298.15 L。
\item 约 5560.54 L，源答案按四位有效数字约 5561 L；实际高压气体非理想性另行复核。
\end{enumerate}
\subsection*{结果与单位}
\noindent 末态体积：\texttt{5560.543} L\par
\noindent\textit{待审核：科学专家审校、版本授权、单题版权、翻译及编排。不可用于正式出卷。}
\medskip\hrule\medskip
\section{给定呼吸耗氧体积的摩尔数}
\noindent\textbf{来源章节：} contents/m51065.md；序号 23。\par
\noindent\textbf{原文链接：} \url{https://github.com/philschatz/chemistry-book/blob/a6ac3918e61cbedf0589c5b43e748dc8a7edf5e2/contents/m51065.md}\par
\noindent\textbf{固定 ID：} \texttt{OSCHEM-de655ee717bf81fa0bd5a579}\par
\textbf{无教材参考答案：} 数值已独立求取，尚待人工复核。\par
\subsection*{独立中文推导}
\begin{enumerate}
\item 给定 14 L、100 kPa 与 25 ℃ 的气态氧，温度必须改为 298.15 K。
\item 使用 R=8.314462618 kPa·L/(mol·K)，代入 n=100×14/(R×298.15)。
\item 得到约 0.56475 mol，按给定精度约 0.56 mol；只针对给定时间的实验模型。
\end{enumerate}
\subsection*{结果与单位}
\noindent 氧气物质的量：\texttt{0.5647536} mol\par
\noindent\textit{待审核：科学专家审校、版本授权、单题版权、翻译及编排。不可用于正式出卷。}
\medskip\hrule\medskip
\section{不同含氟化合物的标准状况体积}
\noindent\textbf{来源章节：} contents/m51065.md；序号 26。\par
\noindent\textbf{原文链接：} \url{https://github.com/philschatz/chemistry-book/blob/a6ac3918e61cbedf0589c5b43e748dc8a7edf5e2/contents/m51065.md}\par
\noindent\textbf{固定 ID：} \texttt{OSCHEM-1b9e55df0faa6023385fd1f7}\par
\subsection*{独立中文推导}
\begin{enumerate}
\item 明确使用教材 STP=273.15 K 和 1 atm；若采用 1 bar 版本 STP 会略有不同。
\item 分别求两种气体的摩尔质量，再用 V=(10.0/M)×RT/P 计算相同质量时的体积。
\item 约 1.85 L 和 4.66 L；分子量越大相同质量的气体物质的量越小。
\end{enumerate}
\subsection*{结果与单位}
\noindent CCl₂F₂：\texttt{1.853818} L\par
\noindent C₂H₅F：\texttt{4.663746} L\par
\noindent\textit{待审核：科学专家审校、版本授权、单题版权、翻译及编排。不可用于正式出卷。}
\medskip\hrule\medskip
\backmatter
\chapter*{人工质量核验清单}
每题均须检查来源版本、化学式、假设与边界、量纲、有效数字、异常答案标签、第三方授权。机器反算不能替代专家审核。
\end{document}
